\documentclass[11pt]{report}

% ============================================================
% PROJECT DOCUMENTATION TEMPLATE
% Structure based on the Sales Assist sample documentation.
% Guidance boxes flag the most common Iteration 1 pitfalls
% (see pitfalls.pdf) at the point in the document where teams
% most often make them. Delete the guidance boxes once you are
% confident a section is complete -- do not leave them in your
% final submission.
% ============================================================

\usepackage[margin=1in]{geometry}
\usepackage{titlesec}
\usepackage{fancyhdr}
\usepackage{enumitem}
\usepackage{longtable}
\usepackage{array}
\usepackage{graphicx}
\usepackage{tcolorbox}
\usepackage{hyperref}
\usepackage{xcolor}

\tcbuselibrary{skins,breakable}

% ---------- Guidance box (remove before final submission) ----------
\newtcolorbox{guidance}[1][]{
  breakable,
  colback=yellow!8,
  colframe=orange!70!black,
  fonttitle=\bfseries,
  title={Watch out: #1},
  boxrule=0.6pt,
  arc=2pt,
  left=6pt, right=6pt, top=4pt, bottom=4pt
}

% ---------- Fill-in placeholder ----------
\newcommand{\fillin}[1]{\textcolor{gray!70}{\emph{[#1]}}}

\pagestyle{fancy}
\fancyhf{}
\rhead{\thepage}
\lhead{\leftmark}
\renewcommand{\headrulewidth}{0.4pt}

\titleformat{\chapter}[hang]{\normalfont\huge\bfseries}{\thechapter.}{1em}{}
\setlist[itemize]{leftmargin=1.5em}

% ============================================================
\title{\textbf{\fillin{Project Name}\\Documentation}}
\author{\fillin{Team Number}\\
\fillin{Member A}\\ \fillin{Member B}\\ \fillin{Member C}\\ \fillin{Member D}\\ \fillin{Member E}}
\date{\fillin{Date}}

\begin{document}

\maketitle
\tableofcontents

% ============================================================
\chapter*{Change Log}
\addcontentsline{toc}{chapter}{Change Log}

\begin{guidance}[Treating early iterations as "draft"]
Do not use placeholder text or "to be decided" sections anywhere in this
document. Each iteration should be written as a complete, reviewable
deliverable -- assume it is final unless a later iteration explicitly
revises it.
\end{guidance}

\textbf{Iteration 1:}
\begin{itemize}
  \item \fillin{What was created/changed in Iteration 1}
\end{itemize}

\textbf{Iteration 2:}
\begin{itemize}
  \item \fillin{What was created/changed in Iteration 2}
\end{itemize}

\textbf{Iteration 3:}
\begin{itemize}
  \item \fillin{What was created/changed in Iteration 3}
\end{itemize}

\textbf{Final Iteration:}
\begin{itemize}
  \item \fillin{What was created/changed in the Final Iteration}
\end{itemize}

% ============================================================
\chapter{Requirements Analysis}

\begin{guidance}[Analysis before implementation]
Finish this section and the Use Cases section \emph{before} you discuss
any classes, code, or UI screens elsewhere in the document. Every
planned feature must trace to at least one requirement and one use case.
Avoid phrases like "we plan to implement \ldots" without first explaining
\emph{why} the feature is needed from the user's perspective.
\end{guidance}

\textbf{Customer / Stakeholder Perspective:}
\fillin{Describe who the end users are and what they are trying to
achieve, in their own words.}

\textbf{Functional Requirements:}
\begin{itemize}
  \item \fillin{Functional requirement 1}
  \item \fillin{Functional requirement 2}
  \item \fillin{Functional requirement 3}
\end{itemize}

\textbf{Non-Functional Requirements / Environment:}
\begin{itemize}
  \item \fillin{Platform, performance, deployment constraints, etc.}
\end{itemize}

% ============================================================
\chapter{User Story}

\begin{guidance}[Requirements traceability]
Every requirement above should be reflected in at least one user story,
and every user story should later map to a use case. Keep IDs or
consistent wording so a reader can trace a requirement all the way
through to the System Sequence Diagrams.
\end{guidance}

\begin{itemize}
  \item As a \fillin{role}, I want to \fillin{goal} so that \fillin{benefit}.
  \item As a \fillin{role}, I want to \fillin{goal} so that \fillin{benefit}.
  \item As a \fillin{role}, I want to \fillin{goal} so that \fillin{benefit}.
\end{itemize}

% ============================================================
\chapter{Key Technical Challenges}

\begin{itemize}
  \item \fillin{Challenge 1 and how it will be approached}
  \item \fillin{Challenge 2 and how it will be approached}
\end{itemize}

\section*{Key Risks}
\begin{itemize}
  \item \fillin{Risk 1}
  \item \fillin{Risk 2}
\end{itemize}

\section*{Competition}
\begin{itemize}
  \item \fillin{Competitor / existing solution 1 -- what it does, what it lacks}
  \item \fillin{Competitor / existing solution 2}
\end{itemize}

% ============================================================
\chapter{What's New?}

\fillin{Summarize what is novel about this project compared to
existing solutions listed under Competition.}

% ============================================================
\chapter{Workflow}

\begin{guidance}[System boundary]
Make sure the workflow diagram clearly separates what the \emph{actor}
does from what the \emph{system} does. Actors should never be shown
performing system responsibilities (e.g., "system authenticates," not
"user authenticates"), and external systems should not be drawn as if
they were internal components.
\end{guidance}

\fillin{Insert workflow / activity diagram here, e.g.\\
\texttt{\textbackslash includegraphics[width=\textbackslash textwidth]\{workflow.png\}}}

% ============================================================
\chapter{Use Cases}

\begin{guidance}[Use cases are not UI click-throughs]
Do not write use cases as a sequence of UI actions ("user clicks
button, system shows page"). Describe the \emph{system's
responsibilities and guarantees} -- business rules, alternative flows,
and constraints -- not screen navigation. UI details will change over
time; system behavior should be the stable part of the documentation.
\end{guidance}

\section{\fillin{Subsystem / Module Name (e.g., Client Application)}}

\textbf{Actor:} \fillin{Actor name}

\textbf{Use Cases:} \fillin{Use case 1}, \fillin{Use case 2}, \fillin{Use case 3}

\textbf{Actor Description:} \fillin{One or two sentences describing the actor.}

\subsection*{Use Case: \fillin{Name}}
\fillin{Short description of the use case.}

\begin{longtable}{|p{0.48\textwidth}|p{0.48\textwidth}|}
\hline
\textbf{Actor: \fillin{Actor}} & \textbf{System: \fillin{System Name}} \\
\hline
1. \fillin{Actor step} & \\
\hline
 & 2. \fillin{System step} \\
\hline
3. \fillin{Actor step} & \\
\hline
\end{longtable}

\bigskip
\textit{Repeat the subsection above for each use case, and repeat the
section above for each subsystem/actor group (e.g., a Server /
Manager-facing module, following the pattern of the Client module).}

% ============================================================
\chapter{Use Case Diagrams}

\fillin{Insert use case diagram(s) here, one per subsystem, e.g.\\
\texttt{\textbackslash includegraphics[width=0.8\textbackslash textwidth]\{use\_case\_diagram.png\}}}

% ============================================================
\chapter{System Sequence Diagram}

\begin{guidance}[Consistency with use cases]
Every scenario shown here should correspond to a use case described
earlier, using the same name. Missing or renamed scenarios break
traceability between Use Cases and SSDs.
\end{guidance}

\section*{\fillin{Scenario Name}}
\fillin{Insert system sequence diagram here.}

\begin{longtable}{|p{0.48\textwidth}|p{0.48\textwidth}|}
\hline
\textbf{Actor} & \textbf{System} \\
\hline
1. \fillin{step} & \\
\hline
 & 2. *\fillin{system response} \\
\hline
3. \fillin{step} & \\
\hline
\end{longtable}

\bigskip
\textit{Repeat for each scenario.}

% ============================================================
\chapter{Domain Model}

\begin{guidance}[Domain model is not a class diagram]
Use only nouns from the problem statement as domain classes. Do not
include methods, data types, UI classes, or persistence/database
details -- those belong to design, not analysis. Show entities and the
associations/multiplicities between them, not implementation.
\end{guidance}

\fillin{Insert domain model diagram here, e.g.\\
\texttt{\textbackslash includegraphics[width=\textbackslash textwidth]\{domain\_model.png\}}}

% ============================================================
\chapter{Test Cases}

\begin{guidance}[Traceability of tests]
Each test case should reference the requirement or use case it
verifies. Keep test case IDs consistent across iterations so gaps in
coverage are easy to spot.
\end{guidance}

\section{Test Cases for Iteration 1}

\textbf{Target Users:} \fillin{Who}

\textbf{What is covered in the tests:}
\begin{enumerate}
  \item \fillin{Area of functionality covered}
\end{enumerate}

\subsection*{Test Case: \fillin{ID, e.g., T1}}
\begin{itemize}
  \item \textbf{Description:} \fillin{What this test checks}
  \item \textbf{Action:} \fillin{Function/action exercised}
  \item \textbf{Related Requirements:} \fillin{Preconditions}
  \item \textbf{Expected Result:} \fillin{Expected behavior}
  \item \textbf{Test Result:} \fillin{Actual behavior observed}
  \item \textbf{Status:} \fillin{Pass / Fail}
\end{itemize}

\bigskip
\textit{Repeat the Test Case block above for each test in this iteration.}

\section{Test Cases for Iteration 2}
\fillin{Same structure as Iteration 1.}

\section{Test Cases for Iteration 3}
\fillin{Same structure as Iteration 1.}

\section{Test Cases for Final Iteration}
\fillin{Same structure as Iteration 1.}

\end{document}